% !PRINTED 345
\weberpage{345}
%<<CONT
gegeben und von \gesperrt{Encke} weiter ausgebildet ist\footnote{Crelle's Journal, Bd.~22 (1841).}, und das sich besonders zu einer praktischen Durchführung der numerischen Rechnungen eignet.

Sind, wie oben, $\alpha, \beta, \gamma \ldots$ die Wurzeln einer Gleichung und ist $\alpha$ die absolut grösste unter ihnen, so dass keine der anderen der Wurzel $\alpha$ an absolutem Werth gleichkommt, so hat
\weq{1}{\sqrt[m]{\alpha^m + \beta^m + \gamma^m + \cdots}}
mit unbegrenzt wachsendem $m$ den Werth $\alpha$ zur Grenze; es ist sogar die Convergenz gegen $\alpha$ eine noch bessere, als bei dem im vorigen Paragraphen betrachteten Quotienten
\[
\frac{\alpha^m + \beta^m + \gamma^m + \cdots}{\alpha^{m-1} + \beta^{m-1} + \gamma^{m-1} + \cdots},
\]
wie man erkennt, wenn man nach dem verallgemeinerten binomischen Lehrsatz:
\[
\begin{aligned}
\sqrt[m]{\alpha^m + \beta^m + \gamma^m + \cdots} &= \alpha\sqrt[m]{1 + \left(\frac{\beta}{\alpha}\right)^m + \left(\frac{\gamma}{\alpha}\right)^m + \cdots} \\
&= \alpha\left[1 + \frac{1}{m}\left(\frac{\beta}{\alpha}\right)^m + \frac{1}{m}\left(\frac{\gamma}{\alpha}\right)^m + \cdots\right]
\end{aligned}
\]
setzt.

Die absolut kleinste unter den Wurzeln erhält man nach demselben Princip als Grenzwerth von
\weq{2}{\dfrac{1}{\sqrt[m]{\dfrac{1}{\alpha^m} + \dfrac{1}{\beta^m} + \dfrac{1}{\gamma^m} + \cdots}},}
und wenn $p$ ein beliebiger Werth ist, so erhält man die dem Werth $p$ am nächsten liegende Wurzel als Grenzwerth von
\weq{3}{p + \dfrac{1}{\sqrt[m]{\dfrac{1}{(\alpha-p)^m} + \dfrac{1}{(\beta-p)^m} + \cdots}}.}
Wenn man die Gleichung $f(x)$ durch die Substitution
\[
y = \frac{1}{x}, \qquad y = \frac{1}{x - p}
\]
transformirt, so gehen die Ausdrücke (2), (3) aus dem Ausdruck (1) hervor, auf den wir uns daher jetzt beschränken wollen.
